\section{Learning without pairs in a perturbation screen}\label{sec:pert}
This note describes the prespecified test of whether intervention-induced population
variation lets RNA--protein dependence within conditions be recovered without
paired cells.

\subsection*{Data and split}
The Perturb-CITE-seq screen of Frangieh et al.\cite{frangieh2021} was obtained
from the scPerturb release\cite{peidli2024} (Zenodo record 7041849): \PtCells{} patient-derived melanoma cells
with CRISPR knockouts of \PtTargets{} genes, cultured alone, with interferon-$\gamma$
or with autologous tumour-infiltrating lymphocytes, with RNA and 24
antibody-derived tags per cell. The release has no replicate or batch
annotation. Targets were read from the detected guides, because the release's
perturbation column names only the first of several targeted genes and labels
cells with targeting and non-targeting guides as controls. We used the
\PtSingleCells{} cells whose guides all targeted one gene and the \PtNTCells{}
cells whose guides were all non-targeting.

Every fourth target in a fixed pseudo-random order was held out (\PtHeldTargets{} targets), and
every cell whose guides targeted a held-out gene was excluded from training and
tuning. Within each training population (target by condition), cells in a fixed
pseudo-random order formed an RNA half and a protein half. Within each
held-out group and each condition of the non-targeting cells, the first half
formed the adaptation cells and the second half the scoring cells, which
alternated between sub-halves A and B. The gene panel was fixed from training
cells only: the \PtEncodingGenes{} genes that encode proteins targeted by the
antibodies and were detected in at least 1\% of training cells, then the genes of highest dispersion
(detected in at least 5\%, not mitochondrial or ribosomal, log dispersion
z-scored within 20 bins of mean log expression) up to \PtGenes{} genes. Protein
features were the centred log ratios of the \PtProteins{} target antibodies. The
split files were recorded before any statistic of held-out or
non-targeting cells other than library sizes was computed.

\subsection*{Protocol}
The protocol was written after a development analysis on
training targets only, in which every fourth training target played the held-out
role (47 targets, 118 groups). There, the channel fitted to
uncentred population means supported two between-condition directions and
recovered \PtDevPooled\% of within-condition dependence, the condition-centred
channel \PtDevPrimary\%, and the paired closed form \PtDevPaired\%; the
condition-centred channel was chosen as primary, because the target is the dependence within
conditions, and the code was run end to end on this split. The protocol, the
code and the training fits were then registered. Predictions for
eligible groups (at least 40 adaptation cells) were computed from adaptation
cells, each assay separately, and recorded
before the scoring files were opened; the evaluation recomputed every
prediction and checked it against the record. There were no amendments.

Arms were the primary condition-centred pairing-free channel with measured RNA
correlation; its matched controls on pseudo-populations and derangements (ten
draws each); the paired closed form, reference regression and pooled
within-population cross-correlation of the condition, fitted on paired training
cells; independence; and, as secondary arms, the uncentred channel with its own
controls, the primary channel with noise-corrected RNA correlation, with its
regression form and with the pooled training correlation matrices of the
condition instead of the group's, the paired closed form with corrected RNA
correlation, and ridge regression on the group's own adaptation pairs. The
recovered fraction (Methods) removes scoring noise from both the gain over
independence and the size of the target, so that arms are compared on the true
within-group dependence; the mean relative Frobenius error, which scoring noise
inflates towards one for every arm, is reported alongside.

\subsection*{Results}
All \PtGroups{} eligible held-out groups came from \PtHeldEligible{} of the
\PtHeldTargets{} held-out targets. The primary channel supported \PtDimS{}
direction (the uncentred channel \PtDimSPooled), and the direction held a median
of \PtCoverage\% of the groups' corrected RNA variance (\PtCoverageNt\% for
non-targeting cells). H1 (recovery above zero) was \PtHOne, H2 (recovery above
both controls) was \PtHTwo, H3 (recovery in non-targeting cells) was \PtHThree{}
and H4 (at least half of the paired closed form's recovery) was \PtHFour. All arms
are in Supplementary Table~\ref{tab:pert}.

\begin{table}[h]
\caption{\textbf{Recovered fraction of within-condition dependence (\%).} Held-out
targets: \PtGroups{} target-by-condition groups, 95\% intervals from 2,000
resamples of targets; conditions are cultured alone, with interferon-$\gamma$ and
with T cells. Non-targeting cells: three conditions pooled, intervals from 2,000
resamples of guide sets. Error, mean relative Frobenius error over held-out
groups against all scoring cells (independence, 1), inflated by scoring noise.
Controls are averages over ten draws.}
\label{tab:pert}
\centering\scriptsize\setlength{\tabcolsep}{3pt}
% Generated by extension/perturbation/make_assets.py. Do not edit.
\begin{tabular}{lcccccc}
\toprule
& \multicolumn{4}{c}{Held-out targets} & \multicolumn{2}{c}{Non-targeting cells} \\
\cmidrule(lr){2-5}\cmidrule(lr){6-7}
Arm & All (95\% CI) & Alone & IFN$\gamma$ & T cells & All (95\% CI) & Error \\
\midrule
No pairs: perturbations, measured RNA (primary) & 6.8 (4.1 to 9.4) & 11.5 & 16.2 & \ensuremath{-}5.4 & 18.6 (16.1 to 21.1) & 0.992 \\
No pairs: perturbations, corrected RNA & 7.0 (4.4 to 9.5) & 11.6 & 16.1 & \ensuremath{-}4.8 & 17.9 (15.5 to 20.3) & 0.991 \\
No pairs: perturbations, regression form & 9.2 (7.9 to 10.5) & 6.0 & 13.2 & 8.0 & -- & 0.989 \\
No pairs: perturbations, training marginals & 15.1 (13.0 to 17.2) & 15.2 & 24.2 & 6.7 & 18.9 (16.5 to 21.3) & 0.982 \\
No pairs: pseudo-populations (10 draws) & \ensuremath{-}0.1 (\ensuremath{-}0.1 to \ensuremath{-}0.1) & -- & -- & -- & -- & -- \\
No pairs: derangements (10 draws) & \ensuremath{-}0.4 (\ensuremath{-}0.6 to \ensuremath{-}0.3) & -- & -- & -- & -- & -- \\
No pairs: conditions pooled & \ensuremath{-}36.4 (\ensuremath{-}42.0 to \ensuremath{-}31.6) & \ensuremath{-}19.7 & \ensuremath{-}26.6 & \ensuremath{-}58.3 & \ensuremath{-}11.9 (\ensuremath{-}16.1 to \ensuremath{-}7.7) & 1.039 \\
No pairs: conditions pooled, pseudo-populations & \ensuremath{-}46.4 (\ensuremath{-}52.3 to \ensuremath{-}41.3) & -- & -- & -- & -- & -- \\
No pairs: conditions pooled, derangements & \ensuremath{-}75.4 (\ensuremath{-}82.0 to \ensuremath{-}69.8) & -- & -- & -- & -- & -- \\
Paired closed form & 37.3 (31.7 to 42.2) & 23.5 & 42.8 & 43.0 & 94.3 (92.8 to 95.7) & 0.953 \\
Paired closed form, corrected RNA & 39.7 (34.4 to 44.4) & 27.4 & 44.8 & 44.6 & -- & 0.950 \\
Reference regression & 53.4 (49.4 to 57.2) & 43.8 & 54.8 & 59.7 & 91.3 (89.6 to 92.9) & 0.935 \\
Transferred correlation & 96.7 (92.7 to 100.7) & 98.3 & 96.6 & 95.5 & 98.5 (97.0 to 99.8) & 0.884 \\
Own paired adaptation cells & 24.0 (19.8 to 27.7) & 9.5 & 25.3 & 34.1 & 89.7 (88.1 to 91.1) & 0.969 \\
\bottomrule
\end{tabular}

\end{table}

\subsection*{Programs and perturbation choice}
These analyses were
specified after the prespecified evaluation and registered before any of their statistics was computed; the
held-out scoring cells had already been opened for the aggregate endpoint, so the analyses are
exploratory. MSigDB Hallmark gene sets were downloaded from gsea-msigdb.org. The interferon block comprised \PgIfnGenes{}
panel genes and the antibodies to HLA class~I (HLA-A, -B, -C and B2M), PD-L1 and
CD47, whose genes belong to the interferon sets. Prespecified criteria: P1, a
pairing-free share of the paired closed form above one half in the interferon
block, was \PgPOne; P2, the same along the supported direction, was \PgPTwo{}
(share \PgAlongShare\%; \PgAlongShareLo{} to \PgAlongShareHi\%); P3, higher
recovery in the interferon block than in the complement (difference \PgDiff\%;
\PgDiffLo--\PgDiffHi\%), was \PgPThree. The support of the interferon program, the
mean squared correlation of its genes with the RNA score along the supported
direction in adaptation cells, was \PgSupportIfn, against at most
\PgSupportOtherMax{} for the \PgOtherCount{} other programs (Supplementary
Table~\ref{tab:programs}).

For perturbation choice, development on training targets only
showed that channels fitted on 10 to 40 targets, selected at
random or by exposure, retained many spurious directions and were often worse
than independence, whereas removing a few targets from the full set left the
estimator stable; the prespecified analysis therefore removed targets. The four
targets of highest exposure, \DsTopFour, had exposures at least \DsExposureGap{}
times that of any other target. After removal of the 5, 10 and 20 targets of
highest exposure, the channel supported \DsFiveDimExp, \DsTenDimExp{} and
\DsTwentyDimExp{} directions, against a mean of \DsFiveDimRand, \DsTenDimRand{} and
\DsTwentyDimRand{} after random removals, none of which removed all. Recovery over
all genes and proteins after removing the 5 of highest exposure was \DsFiveAllExp\%
(random removals, mean \DsFiveAllRand\%, 5th percentile \DsFiveAllRandLow\%), and
in the interferon block \DsFiveIfnExp\% (random \DsFiveIfnRand\%); after removing
10, \DsTenAllExp\% (random \DsTenAllRand\%) and \DsTenIfnExp\% (random
\DsTenIfnRand\%; \DsTenIfnBelow{} of \DsDraws{} random removals were at or below
it). Removing the 5 and 10 targets of strongest effect (\DsTopStrengthFive{} first)
left interferon-block recovery at \DsFiveIfnStr\% and \DsTenIfnStr\%
(random minus strongest, \DsFiveIfnDiffStr\% and \DsTenIfnDiffStr\%; 95\% CI,
\DsFiveIfnDiffStrLo--\DsFiveIfnDiffStrHi\% and \DsTenIfnDiffStrLo--\DsTenIfnDiffStrHi\%).
Effect strength was a comparator, not a criterion. The design criterion was \DsDOne. Leaving single targets out, the loss of
recovery was largest for \DsLooTop, but its rank correlation with exposure over all
targets was \DsLooSpearman{} (permutation $P=\DsLooP$), because targets that move
the same direction substitute for each other.

\begin{table}[h]
\caption{\textbf{Recovery by program (\%).} Held-out target-by-condition groups of
the melanoma screen and non-targeting cells. Blocks: interferon genes and
antibodies; the RNA score along the supported direction with all antibodies;
interferon genes with all antibodies; all other genes and antibodies; genes of
other Hallmark programs with all antibodies. Arms: no pairs, the primary pairing-free
channel; paired CF, paired closed form; ref.\ regr., reference regression; transf.\
corr., pooled within-condition correlation of paired training cells; own pairs, ridge
regression on the group's paired adaptation cells; replicate, the adaptation half
used directly as the estimate. Rel., split-half reliability of the scoring target
(pooled over groups).}
\label{tab:programs}
\centering\scriptsize\setlength{\tabcolsep}{3pt}
% Generated by extension/perturbation/make_assets.py. Do not edit.
\begin{tabular}{lccccccc}
\toprule
Block & No pairs & Paired CF & Ref.\ regr. & Transf.\ corr. & Own pairs & Replicate & Rel.\ \\
\midrule
Interferon block & 42.5 & 64.2 & 68.8 & 94.0 & 48.3 & \ensuremath{-}24.3 & 0.29 \\
Along $S$ & 27.3 & 69.3 & 76.3 & 87.0 & 47.4 & 6.6 & 0.35 \\
Interferon genes, all proteins & 16.4 & 33.5 & 49.1 & 92.3 & 23.9 & \ensuremath{-}281.3 & 0.12 \\
Complement & 5.7 & 36.2 & 54.7 & 97.9 & 18.7 & \ensuremath{-}308.3 & 0.11 \\
E2F targets & \ensuremath{-}5.1 & \ensuremath{-}18.6 & 4.6 & 100.6 & \ensuremath{-}43.9 & \ensuremath{-}554.5 & 0.07 \\
Epithelial--mesenchymal transition & 7.8 & 56.1 & 68.7 & 96.3 & 38.7 & \ensuremath{-}123.2 & 0.18 \\
G2M checkpoint & \ensuremath{-}7.5 & \ensuremath{-}9.4 & 12.3 & 101.3 & \ensuremath{-}43.8 & \ensuremath{-}525.2 & 0.07 \\
Hypoxia & 3.9 & 32.9 & 50.2 & 94.3 & 24.2 & \ensuremath{-}275.1 & 0.12 \\
TNF-$\alpha$ signalling via NF-$\kappa$B & 5.4 & 41.0 & 57.4 & 99.1 & 29.0 & \ensuremath{-}238.2 & 0.13 \\
\midrule
Non-targeting: interferon block & 61.3 & 92.7 & 90.8 & 97.9 & 89.3 & 94.2 & 0.95 \\
Non-targeting: complement & 12.4 & 95.3 & 92.3 & 98.6 & 89.2 & 87.5 & 0.80 \\
\bottomrule
\end{tabular}

\end{table}

\subsection*{Limitations of this test}
The screen used one patient-derived model without biological replicates, so
intervals reflect variation among targets and guides only. Held-out groups had a
median of \PtAdaptMedian{} adaptation cells, which limits the accuracy of the
within-assay covariance matrices that closed-form transfer uses and favours
arms that do not use them. The panel had 20 surface proteins.
